--- Page 1 ---
 
S1 
Supplementary Information for 
 
Quantitative prediction of rate constants and its 
application to organic emitters 
 
Katsuyuki Shizu1 and Hironori Kaji1* 
Email: kaji@scl.kyoto-u.ac.jp 
 
1Institute for Chemical Research, Kyoto University, Uji, Kyoto 611-0011, Japan 
 
 


--- Page 2 ---
 
S2 
Supplementary Method 1 
Method of calculating rate constants 
In this study, each type of rate constant was calculated using the following equations.1,2 
1. S1→S0 fluorescence, kF(S1→S0) 
 
𝑘F(S1 → S0) =
{𝐸(S1) − 𝐸(S0)}3
3𝜋𝜀0ℏ4𝑐3
𝜇(S1-S0)2 × 𝑛ref(𝑛ref
2 + 2)2
9
 
(A1) 
E(S1): energy level of S1 
E(S0): energy level of S0 
μ(S1-S0): transition electric dipole moment between S1 and S0  
nref: refractive index of the surrounding medium (nref = 1.8) 
ε0: vacuum permittivity 
ℏ: Dirac constant (Planck constant divided by 2π) 
c: speed of light 
 
2. Tn→S0 phosphorescence, kPhos(Tn→S0) (n, m ≥ 1) 
 
𝑘Phos(T𝑛 → S0) =
{𝐸(T𝑛) − 𝐸(S0)}3
3𝜋𝜀0ℏ4𝑐3
𝜇(T𝑛-S0)2 × 𝑛ref(𝑛ref
2 + 2)2
9
 
(A2) 
 
𝜇(T𝑛-S0) = {𝜇(T𝑛) − 𝜇(S0)}
SOC(S0-T𝑛)
𝐸(T𝑛) − 𝐸(S0) 
 
 
+ ∑ 𝜇(S𝑚-S0)
SOC(S𝑚-T𝑛)
𝐸(T𝑛) − 𝐸(S𝑚)
𝑚
+ ∑ 𝜇(T𝑛-T𝑚)
SOC(S0-T𝑚)
𝐸(S0) − 𝐸(T𝑚)
𝑚
 
(A3) 
E(Tn): energy level of Tn 
E(Tm): energy level of Tm 
E(Sm): energy level of Sm 
μ(Tn-S0): transition electric dipole moment between Tn and S0 
μ(Tn): permanent electric dipole moment of Tn 
μ(S0): permanent electric dipole moment of S0 
μ(Sm-S0): transition electric dipole moment of Sm 
μ(Tn-Tm): transition electric dipole moment between Tn and Tm 
SOC(S0-Tn): S0-Tn SOC 
SOC(Sm-Tn): S0-Tn SOC 
SOC(S0-Tm): S0-Tm SOC 
 
3. Energy-downhill S1→Tn intersystem crossing, kISC(S1→Tn), and energy-uphill Tn→S1 
reverse intersystem crossing, kRISC(Tn→S1), where E(S1) > E(Tn) 
 
𝑘ISC(S1 → T𝑛) = 2𝜋
ℏ |SOC(S1-T𝑛)|2 × LSFISC(𝐸(S1) − 𝐸(T𝑛)) 
(A4) 


--- Page 3 ---
 
S3 
 
LSFISC(𝐸(S1) − 𝐸(T𝑛)) = 1
𝜋
𝛾
{𝐸(S1) − 𝐸(T𝑛)}2 + 𝛾2 
(A5) 
 
𝑘RISC(T𝑛 → S1) = 1
3 × 𝑘ISC(S1 → T𝑛)exp[−𝛽{𝐸(S1) − 𝐸(T𝑛)}] 
(A6) 
|SOC(S1-Tn)|: norm of the S1-Tn SOC 
γ: half width of half maximum of the line-shape function, LSFISC  
β: inverse temperature (β−1 = kBT) 
kB: Boltzmann constant 
T: temperature 
We used FWHM = 1000 cm-1 (γ = 500 cm−1) according to the previous report.1 
 
4. Energy-uphill S1→Tn intersystem crossing, kISC(S1→Tn), and energy-downhill Tn→S1 
reverse intersystem crossing, kRISC(Tn→S1), where E(S1) < E(Tn) 
 
𝑘RISC(T𝑛 → S1) = 1
3 × 2𝜋
ℏ |SOC(S1-T𝑛)|2 × LSFISC(𝐸(T𝑛) − 𝐸(S1)) 
(A7) 
 
𝑘ISC(S1 → T𝑛) = 3𝑘RISC(T𝑛 → S1)exp[−𝛽{𝐸(T𝑛) − 𝐸(S1)}] 
(A8) 
 
5. S1→S0 nonradiative decay, kNR(S1→S0) 
 
𝑘NR(S1 → S0) = ∑ 𝑘NR,𝛼(S1 → S0)
𝛼
 
(A9) 
 
𝑘NR,𝛼(S1 → S0) = 𝑅𝛼(S1-S0)
ℏ
LSFNR,𝛼 
(A10) 
 
𝑅𝛼(S1-S0) =
ℏ2𝑉𝛼(S1-S0)2
{𝐸(S1) − 𝐸(S0)}2 
(A11) 
 
LSFNR,𝛼 = 2𝜋sinh (1
2 𝛽ℏ𝜔𝛼) ∑ exp {− (1
2 + 𝑣𝛼) 𝛽ℏ𝜔𝛼}
𝑣𝛼
 
 
 
× 1
𝜋𝑆 {
𝑣𝛼𝛾
{𝐸(S1) − 𝐸(S0) + ℏ𝜔𝛼}2 + 𝛾2 +
(𝑣𝛼 + 1)𝛾
{𝐸(S1) − 𝐸(S0) − ℏ𝜔𝛼}2 + 𝛾2} 
(A12) 
Vα(S1-S0): vibronic coupling between S1 and S0 for the αth vibrational mode 
να: vibrational quantum number for the αth vibrational mode 
ωα: angular frequency for the αth vibrational mode 
S: empirical parameter (= 2.1947×105) 
 
6. Energy-downhill Tm→Tn internal conversion, kIC(Tm→Tn), and energy-uphill Tn→Tm 
internal conversion, kIC(Tn→Tm), where E(Tm) > E(Tn) 


--- Page 4 ---
 
S4 
 
𝑘IC(T𝑚 → T𝑛) = ∑ 𝑘IC,𝛼(T𝑚 → T𝑛)
𝛼
 
(A13) 
 
𝑘IC,𝛼(T𝑚 → T𝑛) = 𝑅𝛼(T𝑚-T𝑛)
ℏ
LSFIC,𝛼 
(A14) 
 
𝑅𝛼(T𝑚-T𝑛) =
ℏ2𝑉𝛼(T𝑚-T𝑛)2
{𝐸(T𝑚) − 𝐸(T𝑛)}2 
(A15) 
 
LSFIC,𝛼 = 2𝜋sinh (1
2 𝛽ℏ𝜔𝛼) ∑ exp {− (1
2 + 𝑣𝛼) 𝛽ℏ𝜔𝛼}
𝑣𝛼
 
 
 
× 1
𝜋𝑆 {
𝑣𝛼𝛾
{𝐸(T𝑚) − 𝐸(T𝑛) + ℏ𝜔𝛼}2 + 𝛾2 +
(𝑣𝛼 + 1)𝛾
{𝐸(T𝑚) − 𝐸(T𝑛) − ℏ𝜔𝛼}2 + 𝛾2} 
(A16) 
 
𝑘IC(T𝑛 → T𝑚) = 𝑘IC(T𝑚 → T𝑛)exp[−𝛽{𝐸(T𝑚) − 𝐸(T𝑛)}] 
(A17) 
Vα(Tm-Tn): vibronic coupling between Tm and Tn 
 
7. T1→S0 nonradiative decay, kNR(T1→S0) 
 
𝑘NR(T1 → S0) = 2𝜋
ℏ |SOC(T1-S0)|2 × 1
3𝑆 LSFISC(𝐸(T1) − 𝐸(S0)) 
(A18) 
|SOC(T1-S0)|: norm of the T1-S0 SOC 
 
 


--- Page 5 ---
 
S5 
Supplementary Method 2 
Method of calculating spin-orbit coupling, vibronic coupling constant, transition dipole 
moment, and permanent dipole moment 
Spin-orbit coupling (SOC), vibronic coupling constant, and T1-T2 transition dipole moment 
were calculated using the following formulae. 
1. Spin-orbit coupling between S1 and Tn, |SOC(S1-T𝑛)| 
 
|SOC(S1-T𝑛)| = √
∑
|⟨ΦS1|ℋSOC|ΦT𝑛
𝑀S⟩|
2
𝑀S=0,±1
 
 
(B1) 
 
⟨ΦS1|ℋSOC|ΦT𝑛
𝑀S⟩ = FSC2
2
∑ ∑ 𝑍𝐴
eff ⟨ΦS1| 𝐫𝑖 − 𝐑𝐴
|𝐫𝑖 − 𝐑𝐴|3 × ∇|ΦT𝑛
𝑀S⟩ ∙ 𝐬𝑖
𝐴
𝑖
 
 
(B2) 
ΦS1: electronic wave function of S1 state 
ΦT𝑛
𝑀S: electronic wave function of T1 state 
FSC: fine-structure constant 
𝑍𝐴
eff : effective nuclear charge of the Ath nucleus 
𝐫𝑖: cartesian coordinates of the ith electron,  
𝐑𝐴: cartesian coordinates of the Ath nucleus, 
si: spin operator of the electron i 
 
ΦS1 and ΦT𝑛
𝑀S were calculated with Gaussian 16 Rev C01 program package.3 The integral 
⟨ΦS1|
𝐫𝑖−𝐑𝐴
|𝐫𝑖−𝐑𝐴|3 × ∇|ΦT𝑛
𝑀S⟩ was calculated with the method proposed by McMurchie and 
Davidson.4 
 
2. Vibronic coupling constant between S1 and S0 for the αth vibrational mode, Vα(S1-S0) 
 
𝑉𝛼(S1-S0) = − ∑ ∑ 𝑍𝐴
𝐞MW,𝐴
(𝛼)
√𝑀𝐴
∙ ⟨ΦS1| 𝐫𝑖 − 𝐑𝐴
|𝐫𝑖 − 𝐑𝐴|3 |ΦS0⟩
𝐴
𝑖
 
 
(B3) 
𝑍𝐴: nuclear charge of the Ath nucleus 
𝑀𝐴: mass of the Ath nucleus 
𝐞MW,𝐴
(𝛼)
: the Ath element of the eigenvector of the αth vibrational mode 
ΦS1: electronic wave function of S1 
ΦS0: electronic wave function of S0 
𝐫𝑖: cartesian coordinates of the ith electron,  
𝐑𝐴: cartesian coordinates of the Ath nucleus 
 
ΦS1 and ΦS0 were calculated with Gaussian 16 Rev C01 program package.3 The integral 
⟨ΦS1|
𝐫𝑖−𝐑𝐴
|𝐫𝑖−𝐑𝐴|3 |ΦS0⟩ was calculated with the method proposed by McMurchie and 
Davidson.4 The formula for Tm-Tn vibronic coupling constant is obtained by replacing S1 
and S0 with Tm and Tn, respectively.  
 
 
 


--- Page 6 ---
 
S6 
3. T1-T2 transition dipole moment, 𝜇(T1-T2) 
 
𝜇(T1-T2) = −⟨ΦT1|𝐱|ΦT2⟩ 
 
(B4) 
ΦT1: electronic wave function of T1 
ΦT2: electronic wave function of T2 
𝐱: point in the three-dimensional space 
 
ΦT1 and ΦT2 were calculated with Gaussian 16 Rev C01 program package.3 The integral 
⟨ΦT1|𝐱|ΦT2⟩ was calculated with the method proposed by McMurchie and Davidson.4 
 
4. S1-S0 transition dipole moment and permanent dipole moment of S0, T1, and T2 were 
calculated with Gaussian 16 Rev C01 program package3. 
 


--- Page 7 ---
 
S7 
Supplementary Method 3 
Method of calculating excited-state populations and the number of emitted photons 
Transient photoluminescence decay curve is numerically calculated by the following kinetic 
equations for S0, S1, T1, and T2 populations. The kinetic equations were solved numerically 
using our own code. 
 
𝑑
𝑑𝑡 [S1](𝑡) = −{𝑘ISC(S1 → T1) + 𝑘ISC(S1 → T2) + 𝑘F(S1 → S0) + 𝑘NR(S1 → S0)}[S1](𝑡) 
 
 
                        +𝑘RISC(T1 → S1)[T1](𝑡) + 𝑘RISC(T2 → S1)[T2](𝑡) 
(C1) 
 
𝑑
𝑑𝑡 [T2](𝑡) = −{𝑘RISC(T2 → S1) + 𝑘IC(T2 → T1) + 𝑘Phos(T2 → S0) + 𝑘NR(T2 → S0)}[T2](𝑡) 
 
 
                        +𝑘ISC(S1 → T2)[S1](𝑡) + 𝑘IC(T1 → T2)[T1](𝑡) 
(C2) 
 
𝑑
𝑑𝑡 [T1](𝑡) = −{𝑘RISC(T1 → S1) + 𝑘IC(T1 → T2) + 𝑘Phos(T1 → S0) + 𝑘NR(T1 → S0)}[T1](𝑡) 
 
 
                         +𝑘ISC(S1 → T1)[S1](𝑡) + 𝑘IC(T2 → T1)[T2](𝑡) 
(C3) 
 
𝑑
𝑑𝑡 [S0](𝑡) = {𝑘F(S1 → S0) + 𝑘NR(S1 → S0)}[S1](𝑡) 
 
 
                        +{𝑘Phos(T2 → S0) + 𝑘NR(T2 → S0)}[T2](𝑡) 
 
 
                        +{𝑘Phos(T1 → S0) + 𝑘NR(T1 → S0)}[T1](𝑡) 
(C4) 
 
𝑁(𝑡) = {𝑘F(S1 → S0)[S1](𝑡) + 𝑘Phos(T1 → S0)[T1](𝑡) + 𝑘Phos(T2 → S0)[T2](𝑡)} 
(C5) 
 
Φ = ∫ 𝑁(𝑡)𝑑𝑡 
(C6) 
 
ΦPhos(T1) = ∫ 𝑘Phos(T1 → S0)[T1](𝑡)𝑑𝑡 
(C7) 
 
ΦPhos(T2) = ∫ 𝑘Phos(T2 → S0)[T2](𝑡)𝑑𝑡 
(C8) 
[S0(t)]: population of S0 at time t 
[S1(t)]: population of S1 at time t 
[T1(t)]: population of T1 at time t 
[T2(t)]: population of T2 at time t 
kISC(S1 → T1): rate constant for the S1 → T1 intersystem crossing 
kISC(S1 → T2): rate constant for the S1 → T2 intersystem crossing 
kF(S1 → S0): rate constant for the S1 → S0 radiative decay 
kNR(S1 → S0): rate constant for the S1 → S0 nonradiative decay 
kRISC(T1 → S1): rate constant for the T1 → S1 intersystem crossing 
kRISC(T2 → S1): rate constant for the T2 → S1 intersystem crossing 
kIC(T2 → T1): rate constant for the T2 → T1 internal conversion 
kPhos(T2 → S0): rate constant for the T2 → S0 radiative decay 
kNR(T2 → S0): rate constant for the T2 → S0 nonradiative decay 
kIC(T1 → T2): rate constant for the T1 → T2 internal conversion 
kPhos(T1 → S0): rate constant for the T1 → S0 radiative decay 
kNR(T1 → S0): rate constant for the T1 → S0 nonradiative decay 


--- Page 8 ---
 
S8 
N(t): the number of photons emitted at time t per unit time 
Φ: the total PLQY 
ΦPhos(T1): the contribution from T1 → S0 radiative decay toΦ 
ΦPhos(T2): the contribution from T2 → S0 radiative decay toΦ 
 
In the main text, we define the lifetime for the total radiative decay as τtoR = 1/ktoR and 
mention that τtoR is a simple extension of the averaged radiative decay time proposed by 
Yersin et al.3 Here, we discuss this point in detail. From Equation 9, τtoR is written as 
𝜏toR =
∑
𝑙≥1[S𝑙]
+ ∑
[T𝑙]
𝑙≥1
∑
𝑛≥1 𝑘F(S𝑛 → S0)[S𝑛]
+ ∑
𝑘Phos(T𝑚 → S0)[T𝑚]
𝑚≥1
 
(C9) 
The lifetimes for Sn→S0 fluorescence (τF(Sn→S0)) and Tm→S0 phosphorescence 
(τPhos(Tm→S0)) are written as τF(Sn→S0) = 1/kF(Sn→S0) and τPhos(Tm→S0) = 1/kPhos(Tm→S0), 
respectively (n, m = 1, 2, 3, …). Hence,  
𝜏toR =
∑
𝑙≥1[S𝑙]
+ ∑
[T𝑙]
𝑙≥1
∑
1
𝑛≥1 𝜏F(S𝑛 → S0) [S𝑛]
+ ∑
1
𝜏Phos(T𝑚 → S0) [T𝑚]
𝑚≥1
 
(C10) 
When the excited singlet and triplet states are thermally equilibrated and E(T1) ≤ E(S1), [Sn] 
and [Tm] can be written in terms of [T1] and Sn-T1 and Tm-T1 energy differences 
[S𝑛] = 1
3 [T1]exp (− 𝐸(S𝑛) − 𝐸(T1)
𝑘B𝑇
) 
(C11) 
[T𝑚] = [T1]exp (− 𝐸(T𝑚) − 𝐸(T1)
𝑘B𝑇
) 
(C12) 
Therefore, 
𝜏toR =
∑
𝑒
−𝐸(S𝑙)−𝐸(T1)
𝑘B𝑇
𝑙≥1
+ ∑
3𝑒
−𝐸(T𝑙)−𝐸(T1)
𝑘B𝑇
𝑙≥1
∑
1
𝜏F(S𝑛 → S0) 𝑒
−𝐸(S𝑛)−𝐸(T1)
𝑘B𝑇
𝑛≥1
+ ∑
3
𝜏Phos(T𝑚 → S0) 𝑒
−𝐸(T𝑚)−𝐸(T1)
𝑘B𝑇
𝑚≥1
 
(C13) 
τtoR is the population-weighted average of τF(Sn→S0) and τPhos(Tm→S0). When only S1 and T1 
are considered, C13 can be written as 
𝜏toR =
𝑒
−𝐸(S1)−𝐸(T1)
𝑘B𝑇
+ 3
1
𝜏F(S1 → S0) 𝑒
−𝐸(S1)−𝐸(T1)
𝑘B𝑇
+
3
𝜏Phos(T1 → S0)
 
(C14) 
Yersin et al. proposed the averaged radiative decay time (τav) under the assumption that S1 
and T1 are thermally equilibrated 
𝜏av =
𝑒
−Δ𝐸(S1−T1)
𝑘B𝑇
+ 1
1
𝜏(S1) 𝑒
−Δ𝐸(S1−T1)
𝑘B𝑇
+
1
𝜏(T1)
 
(C15) 
where τ(S1) and τ(T1) are the decay times for S1 and T1, respectively, and ΔE(S1 − T1) is the 
S1-T1 energy gap. Note that the factor of 3 originating from the three triplet substates does not 
appear explicitly in the Yersin et al.'s formula (Equation C15). Comparing Equations C14 and 


--- Page 9 ---
 
S9 
C15 shows that τtoR defined in this study (Equation C14) is a simple extension of the Yersin 
et al.'s formula (Equation C15). 
 
When E(S1) < E(T1) (in the case of inverted singlet-triplet excited states), [Sn] and [Tm] can 
be written in terms of [S1] and Sn-S1 and Tm-S1 energy differences  
[S𝑛] = [S1]exp (− 𝐸(S𝑛) − 𝐸(S1)
𝑘B𝑇
) 
(C16) 
[T𝑚] = 3[S1]exp (− 𝐸(T𝑚) − 𝐸(S1)
𝑘B𝑇
) 
(C17) 
Therefore, 
𝜏toR =
∑
𝑒
−𝐸(S𝑙)−𝐸(S1)
𝑘B𝑇
𝑙≥1
+ ∑
3𝑒
−𝐸(T𝑙)−𝐸(S1)
𝑘B𝑇
𝑙≥1
∑
1
𝜏F(S𝑛 → S0) 𝑒
−𝐸(S𝑛)−𝐸(S1)
𝑘B𝑇
𝑛≥1
+ ∑
3
𝜏Phos(T𝑚 → S0) 𝑒
−𝐸(T𝑚)−𝐸(S1)
𝑘B𝑇
𝑚≥1
 
(C18) 
When only S1 and T1 are considered, C18 can be written as 
𝜏toR =
1 + 3𝑒
−𝐸(T1)−𝐸(S1)
𝑘B𝑇
1
𝜏F(S1 → S0) +
3
𝜏Phos(T1 → S0) 𝑒
−𝐸(T1)−𝐸(S1)
𝑘B𝑇
 
(C19) 
(C19) and (C14) are essentially the same. 
 
In the main text, we define the total ISC ktoISC as 
𝑘toISC = ∑ (∑ 𝑘ISC(S𝑛 → T𝑚)
𝑚≥1
)
[S𝑛]
∑
𝑙≥1[S𝑙]
𝑛≥1
 
(C20) 
When the excited singlet states are thermally equilibrated, from (C16) and (C20) 
𝑘toISC = ∑ (∑ 𝑘ISC(S𝑛 → T𝑚)
𝑚≥1
)
𝑒
−𝐸(S𝑛)−𝐸(S1)
𝑘B𝑇
∑
𝑒
−𝐸(S𝑙)−𝐸(S1)
𝑘B𝑇
𝑙≥1
𝑛≥1
 
(C21) 
Thus, (C20) contains the effect of thermal activation and deactivation between the excited 
singlet states on ktoISC. Likewise, when the triplet states are thermally equilibrated,  
𝑘toRISC = ∑ (∑ 𝑘RISC(T𝑚 → S𝑛)
𝑛≥1
)
[T𝑚]
∑
𝑙≥1[T𝑙]
𝑚≥1
 
(C22) 
can be written as 
𝑘toRISC = ∑ (∑ 𝑘RISC(T𝑚 → S𝑛)
𝑛≥1
)
𝑒
−𝐸(T𝑚)−𝐸(T1)
𝑘B𝑇
∑
𝑒
−𝐸(T𝑙)−𝐸(T1)
𝑘B𝑇
𝑙≥1
𝑚≥1
 
(C23) 
(C22) contains the effect of thermal activation and deactivation between the triplet states on 
ktoRISC. 


--- Page 10 ---
 
S10 
Supplementary Table 1 | Nuclear coordinates of S1 geometry of BNOO in gas phase 
optimised at the TD-TPSSh/6-31G(d) level of theory. 
Atom 
Element symbol 
x (Å)        
y (Å)       
z (Å)        
1 
C 
-0.030591  
4.836110  
0.422831  
2 
C 
0.357717  
6.065813  
0.949661  
3 
C 
1.290846  
6.110435  
1.985771  
4 
C 
1.859043  
4.922663  
2.462596  
5 
C 
1.469632  
3.692895  
1.937428  
6 
C 
-1.469632  
-3.692895  
1.937428  
7 
C 
-1.859043  
-4.922663  
2.462596  
8 
C 
-1.290846  
-6.110435  
1.985771  
9 
C 
-0.357717  
-6.065813  
0.949661  
10 
C 
0.030591  
-4.836110  
0.422831  
11 
C 
0.477826  
3.623503  
0.940309  
12 
C 
-0.477826  
-3.623503  
0.940309  
13 
C 
-0.889482  
1.530142  
-3.127720  
14 
C 
-1.347847  
2.766185  
-3.586806  
15 
C 
-1.363796  
3.878397  
-2.741981  
16 
C 
-0.912728  
3.712275  
-1.429831  
17 
C 
0.912728  
-3.712275  
-1.429831  
18 
C 
1.363796  
-3.878397  
-2.741981  
19 
C 
1.347847  
-2.766185  
-3.586806  
20 
C 
0.889482  
-1.530142  
-3.127720  
21 
N 
0.000000  
-2.419601  
0.394075  
22 
C 
0.458801  
-2.480292  
-0.959730  
23 
C 
0.426667  
-1.331804  
-1.803832  
24 
B 
0.000000  
0.000000  
-1.167356  
25 
C 
-0.426667  
1.331804  
-1.803832  
26 
C 
-0.458801  
2.480292  
-0.959730  
27 
N 
0.000000  
2.419601  
0.394075  
28 
C 
0.000000  
0.000000  
3.214189  
29 
C 
0.066567  
-1.210015  
2.522821  
30 
C 
-0.003217  
-1.204099  
1.109402  
31 
C 
0.000000  
0.000000  
0.378468  
32 
C 
0.003217  
1.204099  
1.109402  
33 
C 
-0.066567  
1.210015  
2.522821  
34 
O 
-0.928886  
4.837814  
-0.611243  
35 
O 
0.928886  
-4.837814  
-0.611243  
36 
H 
-0.075250  
6.967106  
0.528318  
37 
H 
1.593150  
7.068212  
2.397928  
38 
H 
2.620928  
4.953113  
3.235614  
39 
H 
1.933347  
2.776733  
2.286134  
40 
H 
-1.933347  
-2.776733  
2.286134  
41 
H 
-2.620928  
-4.953113  
3.235614  
42 
H 
-1.593150  
-7.068212  
2.397928  
43 
H 
0.075250  
-6.967106  
0.528318  
44 
H 
-0.913067  
0.677366  
-3.799834  
45 
H 
-1.696034  
2.869136  
-4.611150  
46 
H 
-1.701320  
4.856271  
-3.066955  


--- Page 11 ---
 
S11 
47 
H 
1.701320  
-4.856271  
-3.066955  
48 
H 
1.696034  
-2.869136  
-4.611150  
49 
H 
0.913067  
-0.677366  
-3.799834  
50 
H 
0.000000  
0.000000  
4.299906  
51 
H 
0.179196  
-2.141367  
3.066313  
52 
H 
-0.179196  
2.141367  
3.066313  
 
 
Supplementary Table 2 | Nuclear coordinates of S1 geometry of BNSS in gas phase optimised 
at the TD-TPSSh/6-31G(d) level of theory. 
Atom 
Element symbol 
x (Å)        
y (Å)       
z (Å)        
1 
C 
-4.859001  
-0.658323  
0.010570  
2 
C 
-6.005399  
-1.130224  
0.659873  
3 
C 
-5.895403  
-1.861881  
1.842254  
4 
C 
-4.633090  
-2.080683  
2.404577  
5 
C 
-3.489453  
-1.590501  
1.778118  
6 
C 
3.398366  
-1.802436  
-1.663430  
7 
C 
4.507981  
-2.363325  
-2.275492  
8 
C 
5.803840  
-2.072505  
-1.815485  
9 
C 
5.974023  
-1.193705  
-0.752583  
10 
C 
4.861001  
-0.617440  
-0.122061  
11 
C 
-3.583093  
-0.907156  
0.551974  
12 
C 
3.544378  
-0.934614  
-0.554248  
13 
C 
-1.462727  
3.097165  
-0.990891  
14 
C 
-2.651878  
3.553801  
-1.546917  
15 
C 
-3.747434  
2.693649  
-1.678241  
16 
C 
-3.626169  
1.365493  
-1.237089  
17 
C 
3.667300  
1.430923  
1.154184  
18 
C 
3.760018  
2.751101  
1.641889  
19 
C 
2.642551  
3.571274  
1.558308  
20 
C 
1.451475  
3.083892  
1.002957  
21 
N 
2.428237  
-0.388755  
0.081826  
22 
C 
2.476997  
0.947423  
0.599294  
23 
C 
1.317321  
1.774226  
0.504049  
24 
B 
-0.002314  
1.140859  
-0.015799  
25 
C 
-1.302197  
1.767952  
-0.509910  
26 
C 
-2.447485  
0.914751  
-0.634691  
27 
N 
-2.423094  
-0.420037  
-0.108955  
28 
C 
0.001383  
-3.217863  
0.200497  
29 
C 
1.217762  
-2.515777  
0.241788  
30 
C 
1.195619  
-1.122551  
0.096895  
31 
C 
-0.004397  
-0.395924  
0.019933  
32 
C 
-1.208549  
-1.122182  
0.011671  
33 
C 
-1.203543  
-2.538501  
0.063014  
34 
S 
-4.978848  
0.236814  
-1.512802  
35 
S 
5.116382  
0.424878  
1.265246  
36 
H 
-6.979645  
-0.918807  
0.228830  
37 
H 
-6.789127  
-2.236524  
2.332639  


--- Page 12 ---
 
S12 
38 
H 
-4.537737  
-2.614668  
3.345750  
39 
H 
-2.513936  
-1.732474  
2.231893  
40 
H 
2.401530  
-2.005684  
-2.037414  
41 
H 
4.368765  
-3.015735  
-3.132168  
42 
H 
6.670337  
-2.510137  
-2.301129  
43 
H 
6.969014  
-0.935695  
-0.402386  
44 
H 
-0.607952  
3.766884  
-0.949849  
45 
H 
-2.731999  
4.580432  
-1.895798  
46 
H 
-4.680475  
3.030639  
-2.118164  
47 
H 
4.699755  
3.110360  
2.049138  
48 
H 
2.695771  
4.594343  
1.920204  
49 
H 
0.586091  
3.738744  
0.962949  
50 
H 
0.005789  
-4.301725  
0.263180  
51 
H 
2.153995  
-3.048076  
0.375171  
52 
H 
-2.136325  
-3.088063  
-0.007886  
 
 
Supplementary Table 3 | Nuclear coordinates of S1 geometry of BNSeSe in gas phase 
optimised using the TD-TPSSh method. For the H, B, C, and N atoms, the 6-31G(d) basis set 
was used. For the Se atoms, the Stuttgart/Dresden pseudopotentials and basis set (SDD) were 
used. 
Atom 
Element symbol 
x (Å)        
y (Å)       
z (Å)        
1 
C 
0.405028  
4.837447  
0.756781  
2 
C 
1.171868  
5.879949  
1.285342  
3 
C 
2.338372  
5.604678  
1.997899  
4 
C 
2.760589  
4.277484  
2.150037  
5 
C 
2.012980  
3.237532  
1.611024  
6 
C 
-2.012980  
-3.237532  
1.611024  
7 
C 
-2.760589  
-4.277484  
2.150037  
8 
C 
-2.338372  
-5.604678  
1.997899  
9 
C 
-1.171868  
-5.879949  
1.285342  
10 
C 
-0.405028  
-4.837447  
0.756781  
11 
C 
0.803417  
3.496113  
0.928709  
12 
C 
-0.803417  
-3.496113  
0.928709  
13 
C 
-0.900012  
1.514751  
-3.096470  
14 
C 
-1.403539  
2.728242  
-3.562394  
15 
C 
-1.484990  
3.830002  
-2.709591  
16 
C 
-1.036708  
3.691136  
-1.385359  
17 
C 
1.036708  
-3.691136  
-1.385359  
18 
C 
1.484990  
-3.830002  
-2.709591  
19 
C 
1.403539  
-2.728242  
-3.562394  
20 
C 
0.900012  
-1.514751  
-3.096470  
21 
N 
-0.043534  
-2.423295  
0.409126  
22 
C 
0.489382  
-2.493276  
-0.925481  
23 
C 
0.427074  
-1.338011  
-1.773343  
24 
B 
0.000000  
0.000000  
-1.148036  
25 
C 
-0.427074  
1.338011  
-1.773343  
26 
C 
-0.489382  
2.493276  
-0.925481  


--- Page 13 ---
 
S13 
27 
N 
0.043534  
2.423295  
0.409126  
28 
C 
0.000000  
0.000000  
3.225231  
29 
C 
0.000000  
-1.212354  
2.532475  
30 
C 
-0.035167  
-1.200687  
1.117554  
31 
C 
0.000000  
0.000000  
0.392758  
32 
C 
0.035167  
1.200687  
1.117554  
33 
C 
0.000000  
1.212354  
2.532475  
34 
Se 
-1.242651  
5.206223  
-0.180429  
35 
Se 
1.242651  
-5.206223  
-0.180429  
36 
H 
0.851538  
6.906424  
1.132851  
37 
H 
2.926326  
6.419773  
2.409263  
38 
H 
3.688408  
4.054438  
2.668532  
39 
H 
2.358246  
2.212327  
1.694538  
40 
H 
-2.358246  
-2.212327  
1.694538  
41 
H 
-3.688408  
-4.054438  
2.668532  
42 
H 
-2.926326  
-6.419773  
2.409263  
43 
H 
-0.851538  
-6.906424  
1.132851  
44 
H 
-0.892706  
0.657411  
-3.763249  
45 
H 
-1.745384  
2.818045  
-4.590412  
46 
H 
-1.877652  
4.781246  
-3.053938  
47 
H 
1.877652  
-4.781246  
-3.053938  
48 
H 
1.745384  
-2.818045  
-4.590412  
49 
H 
0.892706  
-0.657411  
-3.763249  
50 
H 
0.000000  
0.000000  
4.310809  
51 
H 
0.036833  
-2.153693  
3.070726  
52 
H 
-0.036833  
2.153693  
3.070726  
 
 
Supplementary Table 4 | Nuclear coordinates of S1 geometry of BNTeTe in gas phase 
optimised using the TD-TPSSh method. For the H, B, C, and N atoms, the 6-31G(d) basis set 
was used. For the Te atoms, the Stuttgart/Dresden pseudopotentials and basis set (SDD) were 
used. 
Atom 
Element symbol 
x (Å)        
y (Å)       
z (Å)        
1 
C 
-4.776822  
0.826524  
-0.835490  
2 
C 
-5.696138  
1.352599  
-1.752086  
3 
C 
-5.258996  
1.958018  
-2.929218  
4 
C 
-3.888258  
2.009073  
-3.214218  
5 
C 
-2.968384  
1.474101  
-2.321646  
6 
C 
2.968384  
1.474101  
2.321646  
7 
C 
3.888258  
2.009073  
3.214218  
8 
C 
5.258996  
1.958018  
2.929218  
9 
C 
5.696138  
1.352599  
1.752086  
10 
C 
4.776822  
0.826524  
0.835490  
11 
C 
-3.392862  
0.893659  
-1.104589  
12 
C 
3.392862  
0.893659  
1.104589  
13 
C 
-1.562406  
-3.112009  
0.812702  
14 
C 
-2.798431  
-3.576708  
1.257576  
15 
C 
-3.899881  
-2.720693  
1.288627  


--- Page 14 ---
 
S14 
16 
C 
-3.745937  
-1.394460  
0.841903  
17 
C 
3.745937  
-1.394460  
-0.841903  
18 
C 
3.899881  
-2.720693  
-1.288627  
19 
C 
2.798431  
-3.576708  
-1.257576  
20 
C 
1.562406  
-3.112009  
-0.812702  
21 
N 
2.416484  
0.386903  
0.209412  
22 
C 
2.523097  
-0.943189  
-0.338207  
23 
C 
1.363346  
-1.790815  
-0.344314  
24 
B 
0.000000  
-1.168132  
0.000000  
25 
C 
-1.363346  
-1.790815  
0.344314  
26 
C 
-2.523097  
-0.943189  
0.338207  
27 
N 
-2.416484  
0.386903  
-0.209412  
28 
C 
0.000000  
3.202138  
0.000000  
29 
C 
1.209190  
2.509183  
0.092723  
30 
C 
1.195228  
1.094334  
0.117614  
31 
C 
0.000000  
0.371032  
0.000000  
32 
C 
-1.195228  
1.094334  
-0.117614  
33 
C 
-1.209190  
2.509183  
-0.092723  
34 
Te 
-5.391290  
-0.063430  
0.974232  
35 
Te 
5.391290  
-0.063430  
-0.974232  
36 
H 
-6.759814  
1.284295  
-1.542137  
37 
H 
-5.982138  
2.366354  
-3.629219  
38 
H 
-3.538141  
2.444558  
-4.145572  
39 
H 
-1.908980  
1.478503  
-2.556374  
40 
H 
1.908980  
1.478503  
2.556374  
41 
H 
3.538141  
2.444558  
4.145572  
42 
H 
5.982138  
2.366354  
3.629219  
43 
H 
6.759814  
1.284295  
1.542137  
44 
H 
-0.704656  
-3.777405  
0.851105  
45 
H 
-2.905137  
-4.603879  
1.597620  
46 
H 
-4.864617  
-3.069533  
1.643438  
47 
H 
4.864617  
-3.069534  
-1.643438  
48 
H 
2.905137  
-4.603879  
-1.597620  
49 
H 
0.704656  
-3.777405  
-0.851105  
50 
H 
0.000000  
4.287743  
0.000000  
51 
H 
2.151111  
3.046410  
0.133236  
52 
H 
-2.151111  
3.046410  
-0.133236  
 
 
Supplementary Table 5 | Nuclear coordinates of S1 geometry of BNPoPo in gas phase 
optimised using the TD-TPSSh method. For the H, B, C, and N atoms, the 6-31G(d) basis set 
was used. For the Po atoms, the Stuttgart/Dresden pseudopotentials and basis set (SDD) were 
used. 
Atom 
Element symbol 
x (Å)        
y (Å)       
z (Å)        
1 
C 
-4.782300  
-0.592113  
1.251140  
2 
C 
-5.663829  
-0.827989  
2.312850  
3 
C 
-5.174519  
-1.073153  
3.596607  
4 
C 
-3.794939  
-1.037867  
3.826902  


--- Page 15 ---
 
S15 
5 
C 
-2.916155  
-0.771558  
2.780102  
6 
C 
2.668356  
-1.990067  
-2.058160  
7 
C 
3.470808  
-2.678496  
-2.949948  
8 
C 
4.872380  
-2.563401  
-2.899216  
9 
C 
5.450151  
-1.737767  
-1.944454  
10 
C 
4.646995  
-1.031230  
-1.033607  
11 
C 
-3.395155  
-0.567250  
1.471304  
12 
C 
3.229186  
-1.138976  
-1.067597  
13 
C 
-1.552199  
2.974020  
-1.150174  
14 
C 
-2.789842  
3.389491  
-1.625511  
15 
C 
-3.903986  
2.549203  
-1.510356  
16 
C 
-3.753379  
1.299247  
-0.888765  
17 
C 
3.781130  
1.440434  
0.552803  
18 
C 
3.953681  
2.807135  
0.856573  
19 
C 
2.847732  
3.645065  
0.786687  
20 
C 
1.596937  
3.118530  
0.441175  
21 
N 
2.375742  
-0.469191  
-0.186097  
22 
C 
2.542210  
0.912239  
0.175652  
23 
C 
1.385923  
1.761988  
0.133888  
24 
B 
-0.012524  
1.121627  
-0.109957  
25 
C 
-1.360444  
1.727105  
-0.494789  
26 
C 
-2.532748  
0.908820  
-0.330911  
27 
N 
-2.457409  
-0.321679  
0.414146  
28 
C 
-0.048022  
-3.122237  
0.838552  
29 
C 
1.170404  
-2.481872  
0.546538  
30 
C 
1.142797  
-1.137959  
0.172762  
31 
C 
-0.041663  
-0.384841  
0.155097  
32 
C 
-1.246050  
-1.046985  
0.458245  
33 
C 
-1.250044  
-2.422377  
0.779792  
34 
Po 
-5.449926  
-0.142606  
-0.815616  
35 
Po 
5.574887  
0.167615  
0.541566  
36 
H 
-6.735869  
-0.821664  
2.135337  
37 
H 
-5.864756  
-1.272077  
4.412018  
38 
H 
-3.402199  
-1.196340  
4.827600  
39 
H 
-1.847237  
-0.718853  
2.962904  
40 
H 
1.588964  
-2.063372  
-2.121876  
41 
H 
3.005070  
-3.297201  
-3.711324  
42 
H 
5.496720  
-3.103559  
-3.604177  
43 
H 
6.531025  
-1.632995  
-1.904470  
44 
H 
-0.684293  
3.607283  
-1.316258  
45 
H 
-2.893164  
4.356534  
-2.112706  
46 
H 
-4.869257  
2.855119  
-1.902351  
47 
H 
4.932140  
3.198146  
1.119794  
48 
H 
2.954744  
4.704136  
1.005916  
49 
H 
0.738935  
3.784119  
0.416283  
50 
H 
-0.050632  
-4.177267  
1.095737  
51 
H 
2.106596  
-3.030494  
0.590767  
52 
H 
-2.189707  
-2.927114  
0.981350  
 


--- Page 16 ---
 
S16 
Supplementary Table 6 | Nuclear coordinates of S1 geometry of BNCOCO in gas phase 
optimised at the TD-TPSSh/6-31G(d) level of theory. 
Atom 
Element symbol 
x (Å)        
y (Å)       
z (Å)        
1 
 C 
-0.048300  
4.859033  
0.570623  
2 
 C 
0.359951  
6.031304  
1.223221  
3 
 C 
1.289809  
5.988591  
2.257187  
4 
 C 
1.855451  
4.758720  
2.624923  
5 
 C 
1.452213  
3.580660  
2.001920  
6 
 C 
-1.452213  
-3.580660  
2.001920  
7 
 C 
-1.855451  
-4.758720  
2.624923  
8 
 C 
-1.289809  
-5.988591  
2.257187  
9 
 C 
-0.359951  
-6.031304  
1.223221  
10 
 C 
0.048300  
-4.859033  
0.570623  
11 
 C 
0.455376  
3.612671  
1.005296  
12 
 C 
-0.455376  
-3.612671  
1.005296  
13 
 C 
-0.936319  
1.474792  
-3.104385  
14 
 C 
-1.423890  
2.689509  
-3.587524  
15 
 C 
-1.428731  
3.812528  
-2.766092  
16 
 C 
-0.934864  
3.731489  
-1.447809  
17 
 C 
0.934864  
-3.731489  
-1.447809  
18 
 C 
1.428731  
-3.812528  
-2.766092  
19 
 C 
1.423890  
-2.689509  
-3.587524  
20 
 C 
0.936319  
-1.474792  
-3.104385  
21 
 N 
0.000000  
-2.424326  
0.386795  
22 
 C 
0.468054  
-2.490506  
-0.960449  
23 
 C 
0.452613  
-1.324365  
-1.779277  
24 
 B 
0.000000  
0.000000  
-1.151837  
25 
 C 
-0.452613  
1.324365  
-1.779277  
26 
 C 
-0.468054  
2.490506  
-0.960449  
27 
 N 
0.000000  
2.424326  
0.386795  
28 
 C 
0.000000  
0.000000  
3.202583  
29 
 C 
0.095612  
-1.208684  
2.514391  
30 
 C 
0.007884  
-1.210699  
1.097326  
31 
 C 
0.000000  
0.000000  
0.375336  
32 
 C 
-0.007884  
1.210699  
1.097326  
33 
 C 
-0.095612  
1.208684  
2.514391  
34 
 H 
-0.063887  
6.966108  
0.869304  
35 
 H 
1.603992  
6.903576  
2.750507  
36 
 H 
2.630820  
4.717631  
3.384281  
37 
 H 
1.931195  
2.641670  
2.257691  
38 
 H 
-1.931195  
-2.641670  
2.257691  
39 
 H 
-2.630820  
-4.717631  
3.384281  
40 
 H 
-1.603992  
-6.903576  
2.750507  
41 
 H 
0.063887  
-6.966108  
0.869304  
42 
 H 
-0.956216  
0.598692  
-3.745895  
43 
 H 
-1.799744  
2.758031  
-4.604993  
44 
 H 
-1.796598  
4.776282  
-3.101753  
45 
 H 
1.796598  
-4.776282  
-3.101753  


--- Page 17 ---
 
S17 
46 
 H 
1.799744  
-2.758031  
-4.604993  
47 
 H 
0.956216  
-0.598692  
-3.745895  
48 
 H 
0.000000  
0.000000  
4.287953  
49 
 H 
0.242303  
-2.133480  
3.059394  
50 
 H 
-0.242303  
2.133480  
3.059394  
51 
 C 
0.905240  
-4.949405  
-0.626994  
52 
 C 
-0.905240  
4.949405  
-0.626994  
53 
 O 
1.457783  
-6.005610  
-0.963549  
54 
 O 
-1.457783  
6.005610  
-0.963549  
 
 


--- Page 18 ---
 
S18 
 
Supplementary Figure 1 | Ball-and-stick representations of the optimised S1 geometries of 
(a) BNOO, (b) BNSS, (c) BNSeSe, (d) BNTeTe, (e) BNPoPo, and (f) BNCOCO. The N-C, 
B-C, and X-C (X = O, S, Se, Te, Po, and CO) bond lengths are in Å (left). The dihedral 
angles are in degree (right).  
 
 


--- Page 19 ---
 
S19 
Supplementary Table 7 | S1, S2, T1, T2, and T3 of BNOO calculated using the TD-
B3LYP//TD-TPSSh method. Φ𝑎
𝑟 is a Slater determinant that denotes the single-electron 
excitation from the ath occupied orbital to the rth unoccupied orbital. 
State 
Excitation 
energy (eV) 
Transition dipole  
moment (au) 
Slater 
determinant 
Contribution (%) 
T1 
2.0902 
 
ΦHOMO
LUMO       
97.55 
T2 
2.4251 
 
ΦHOMO−1
LUMO
 
93.36 
 
 
 
ΦHOMO
LUMO+5 
1.28 
S1 
2.5052 
1.702 
ΦHOMO
LUMO 
98.82 
T3 
2.8443 
 
ΦHOMO
LUMO+1 
79.83 
 
 
 
ΦHOMO
LUMO+7 
3.89 
 
 
 
ΦHOMO−3
LUMO+7 
1.98 
 
 
 
ΦHOMO−1
LUMO+2 
1.72 
 
 
 
ΦHOMO−2
LUMO+1 
1.15 
 
 
 
ΦHOMO−5
LUMO+1 
1.06 
 
 
 
ΦHOMO−4
LUMO
 
1.06 
S2 
2.9575 
0.909 
ΦHOMO−1
LUMO
 
98.58 
 
 
Supplementary Table 8 | S1, S2, T1, T2, and T3 of BNSS calculated using the TD-
B3LYP//TD-TPSSh method. Φ𝑎
𝑟 is a Slater determinant that denotes the single-electron 
excitation from the ath occupied orbital to the rth unoccupied orbital. 
State 
Excitation 
energy (eV) 
Transition dipole  
moment (au) 
Slater 
determinant 
Contribution (%) 
T1 
2.0879 
 
ΦHOMO
LUMO       
96.97 
S1 
2.4388 
1.607 
ΦHOMO
LUMO 
99.10 
T2 
2.4742 
 
ΦHOMO−1
LUMO
 
92.93 
 
 
 
ΦHOMO
LUMO+5 
1.46 
T3 
2.8686 
 
ΦHOMO
LUMO+1 
51.13 
 
 
 
ΦHOMO
LUMO+3 
10.61 
 
 
 
ΦHOMO−2
LUMO
 
6.29 
 
 
 
ΦHOMO
LUMO+2 
5.67 
 
 
 
ΦHOMO−1
LUMO+3 
3.73 
 
 
 
ΦHOMO
LUMO+6 
1.61 
 
 
 
ΦHOMO−5
LUMO+1 
1.23 
 
 
 
ΦHOMO−2
LUMO+1 
1.16 
 
 
 
ΦHOMO−1
LUMO+5 
1.16 


--- Page 20 ---
 
S20 
 
 
 
ΦHOMO−6
LUMO
 
1.04 
S2 
2.9046 
0.892 
ΦHOMO−1
LUMO
 
98.72 
 
 
Supplementary Table 9 | S1, S2, T1, T2, and T3 of BNSeSe calculated using the TD-
B3LYP//TD-TPSSh method. Φ𝑎
𝑟 is a Slater determinant that denotes the single-electron 
excitation from the ath occupied orbital to the rth unoccupied orbital. 
State 
Excitation 
energy (eV) 
Transition dipole  
moment (au) 
Slater 
determinant 
Contribution (%) 
T1 
2.1630 
 
ΦHOMO
LUMO       
97.19 
S1 
2.5105 
1.653 
ΦHOMO
LUMO 
99.03 
T2 
2.5557 
 
ΦHOMO−1
LUMO
 
90.09 
 
 
 
ΦHOMO
LUMO+5 
2.04 
 
 
 
ΦHOMO−3
LUMO
 
1.98 
 
 
 
ΦHOMO−7
LUMO
 
1.18 
T3 
2.9395 
 
ΦHOMO
LUMO+1 
69.04 
 
 
 
ΦHOMO−2
LUMO+1 
2.87 
 
 
 
ΦHOMO
LUMO+5 
2.74 
 
 
 
ΦHOMO−4
LUMO
 
2.60 
 
 
 
ΦHOMO
LUMO+7 
2.52 
 
 
 
ΦHOMO−1
LUMO
 
2.44 
 
 
 
ΦHOMO−1
LUMO+6 
1.76 
 
 
 
ΦHOMO
LUMO+11 
1.72 
 
 
 
ΦHOMO−3
LUMO+6 
1.47 
 
 
 
ΦHOMO−8
LUMO+11 
1.46 
 
 
 
ΦHOMO−5
LUMO+1 
1.11 
S2 
2.9543 
0.498 
ΦHOMO−1
LUMO
 
98.92 
 
 
Supplementary Table 10 | S1, S2, T1, T2, and T3 of BNTeTe calculated using the TD-
B3LYP//TD-TPSSh method. Φ𝑎
𝑟 is a Slater determinant that denotes the single-electron 
excitation from the ath occupied orbital to the rth unoccupied orbital. 
State 
Excitation 
energy (eV) 
Transition dipole  
moment (au) 
Slater 
determinant 
Contribution (%) 
T1 
2.1827 
 
ΦHOMO
LUMO       
96.36 
 
 
 
ΦHOMO−2
LUMO
 
1.34 
S1 
2.5158 
1.6154 
ΦHOMO
LUMO 
99.01 


--- Page 21 ---
 
S21 
T2 
2.5744 
 
ΦHOMO−1
LUMO
 
88.03 
 
 
 
ΦHOMO−3
LUMO
 
4.46 
 
 
 
ΦHOMO
LUMO+5 
1.84 
 
 
 
ΦHOMO−7
LUMO
 
1.21 
T3 
2.8780 
 
ΦHOMO
LUMO+2 
26.51 
 
 
 
ΦHOMO−1
LUMO+1 
24.82 
 
 
 
ΦHOMO−2
LUMO
 
23.49 
 
 
 
ΦHOMO−2
LUMO+2 
8.27 
 
 
 
ΦHOMO−5
LUMO
 
4.27 
 
 
 
ΦHOMO−6
LUMO
 
1.65 
S2 
2.9086 
0.6065 
ΦHOMO−1
LUMO
 
98.44 
 
 
Supplementary Table 11 | S1, S2, T1, T2, and T3 of BNPoPo calculated using the TD-
B3LYP//TD-TPSSh method. Φ𝑎
𝑟 is a Slater determinant that denotes the single-electron 
excitation from the ath occupied orbital to the rth unoccupied orbital. 
State 
Excitation 
energy (eV) 
Transition dipole  
moment (au) 
Slater 
determinant 
Contribution (%) 
T1 
2.1339 
 
ΦHOMO
LUMO       
92.84 
 
 
 
ΦHOMO−2
LUMO
 
2.31 
 
 
 
ΦHOMO−1
LUMO
 
2.15 
S1 
2.4062 
1.3456 
ΦHOMO
LUMO 
99.23 
T2 
2.5265 
 
ΦHOMO−1
LUMO
 
84.24 
 
 
 
ΦHOMO−3
LUMO
 
5.28 
 
 
 
ΦHOMO−2
LUMO+2 
2.90 
 
 
 
ΦHOMO
LUMO 
2.15 
T3 
2.6655 
 
ΦHOMO
LUMO+2 
81.27 
 
 
 
ΦHOMO−1
LUMO+2 
9.07 
 
 
 
ΦHOMO−2
LUMO+2 
4.75 
S2 
2.8286 
0.9483 
ΦHOMO−1
LUMO
 
98.58 
 
 
Supplementary Table 12 | S1, S2, T1, T2, and T3 of BNCOCO calculated using the TD-
B3LYP//TD-TPSSh method. Φ𝑎
𝑟 is a Slater determinant that denotes the single-electron 
excitation from the ath occupied orbital to the rth unoccupied orbital. 
State 
Excitation 
energy (eV) 
Transition dipole  
moment (au) 
Slater 
determinant 
Contribution (%) 


--- Page 22 ---
 
S22 
T1 
2.2005 
 
ΦHOMO
LUMO       
95.56 
 
 
 
ΦHOMO−1
LUMO+1 
2.22 
T2 
2.5191 
 
ΦHOMO
LUMO+1 
51.81 
 
 
 
ΦHOMO−1
LUMO
 
33.11 
 
 
 
ΦHOMO
LUMO+3 
4.78 
 
 
 
ΦHOMO−1
LUMO+2 
1.58 
 
 
 
ΦHOMO−5
LUMO+1 
1.25 
S1 
2.5454 
2.2007 
ΦHOMO
LUMO 
98.61 
T3 
2.8368 
 
ΦHOMO
LUMO+2 
33.91 
 
 
 
ΦHOMO−1
LUMO+1 
26.84 
 
 
 
ΦHOMO−5
LUMO
 
16.86 
 
 
 
ΦHOMO−3
LUMO
 
4.17 
 
 
 
ΦHOMO−11
LUMO+1  
1.83 
 
 
 
ΦHOMO−10
LUMO
 
1.57 
 
 
 
ΦHOMO−4
LUMO+1 
1.11 
 
 
 
ΦHOMO−6
LUMO
 
1.10 
 
 
 
ΦHOMO−2
LUMO+1 
1.08 
 
 
 
ΦHOMO−8
LUMO
 
1.08 
S2 
3.1061 
0.2341 
ΦHOMO
LUMO+1 
91.50 
 
 
 
ΦHOMO−1
LUMO
 
5.47 
 
 
 
 
 


--- Page 23 ---
 
S23 
Supplementary Method 4 
T1-energy correction with the combined TD-B3LYP and TDA-B2-PLYP method 
First, we calculated the S1 and T2 energies (denoted EB3LYP(S1) and EB3LYP(T2)) using the TD-
B3LYP/6-31G(d) (and SDD) method for the optimised S1 geometry. Then, we calculated 
ΔE(T1→S1) (denoted as ΔEB2-PLYP(T1→S1)) using the TDA-B2-PLYP/def2-TZVP method for 
the same geometry. Finally, we calculated the T1 energy as E(T1) = EB3LYP(S1) − ΔEB2-
PLYP(T1→S1) (Table S13). For the S1 and T2 energies, we used the TD-B3LYP/6-31G(d) 
results without correction: E(S1) = EB3LYP(S1) and E(T2) = EB3LYP(T2). We used the E(S1), 
E(T1), and E(T2) values to calculate the rate constants. 
 
 
 


--- Page 24 ---
 
S24 
Supplementary Table 13 | Calculated excited-state energies and energy differences with 
the combined TD-B3LYP and TDA-B2-PLYP method in eV. The TD-B3LYP calculations 
were performed with Gaussian 16 Rev C01 program package.3 The TDA-B2-PLYP 
calculations were performed with ORCA 5.0.3 program package.6-8 
 
BNOO 
BNSS BNSeSe BNTeTe BNPoPo BNCOCO 
E(S1) = EB3LYP(S1) 
2.51 
2.44 
2.51 
2.52 
2.41 
2.55 
E(T2) = EB3LYP(T2) 
2.43 
2.47 
2.56 
2.57 
2.53 
2.52 
ΔEB3LYP(T1→S1) 
0.42 
0.35 
0.35 
0.33 
0.27 
0.35 
ΔE(T1→S1) = ΔEB2-PLYP(T1→S1) 
0.21 
0.14 
0.14 
0.14 
0.14 
0.14 
E(T1) = E(S1) − ΔE(T1→S1) 
2.30 
2.30 
2.37 
2.38 
2.26 
2.41 
ΔE(T2→S1) = E(S1) − E(T2) 
0.08 
−0.04 
−0.05 
−0.06 
−0.12 
0.03 
ΔE(T1→T2) = E(T2) − E(T1) 
0.13 
0.18 
0.19 
0.20 
0.26 
0.11 
Experiments 
 
 
 
 
 
 
E(S1) 
2.54* 
2.55* 
2.58* 
 
 
 
E(T1) 
2.39* 
2.42* 
2.44* 
 
 
 
ΔE(T1→S1) 
0.15* 
0.13* 
0.14* 
 
 
 
*) Experimental data from ref. 9. 
 
 


--- Page 25 ---
 
S25 
Supplementary Table 14 | Calculated excited-state energies and energy differences with 
the ADC(2) and SCS-CC2 methods in eV. The def2-TZVP basis set was used for all the 
atoms. The ADC(2) and SCS-CC2 calculations were performed with TURBOMOLE v7.4.1 
2019 program package.10 
 
BNOO 
BNSS 
BNSeSe 
BNTeTe 
BNPoPo BNCOCO 
ADC(2) 
 
 
 
 
 
 
EADC2(S1) 
2.84 
2.30 
2.35 
2.35 
2.33 
3.07 
EADC2(T1) 
2.61 
2.22 
2.27 
2.27 
2.25 
2.83 
EADC2(T2) 
3.14 
2.67 
3.23 
2.77 
2.71 
3.03 
ΔEADC2(T1→S1) 
0.23 
0.09 
0.08 
0.08 
0.07 
0.24 
ΔEADC2(T2→S1) 
-0.30 
-0.37 
-0.88 
-0.42 
-0.39 
0.04 
ΔEADC2(T1→T2) 
0.53 
0.45 
0.96 
0.49 
0.46 
0.20 
SCS-CC2 
 
 
 
 
 
 
ESCS-CC2(S1) 
2.68 
2.68 
2.72 
2.72 
2.72 
3.36 
ESCS-CC2(T1) 
2.58 
2.58 
2.63 
2.64 
2.64 
3.09 
ESCS-CC2(T2) 
3.02 
3.02 
3.11 
3.12 
2.97 
3.49 
ΔESCS-CC2(T1→S1) 
0.10 
0.10 
0.09 
0.09 
0.08 
0.27 
ΔESCS-CC2(T2→S1) 
-0.34 
-0.34 
-0.39 
-0.39 
-0.25 
-0.13 
ΔESCS-CC2(T1→T2) 
0.43 
0.43 
0.48 
0.48 
0.33 
0.40 
 
 
 
 


--- Page 26 ---
 
S26 
 
 
Supplementary Figure 2 | Calculated population, population ratio, and rate constants for total 
radiative decay for BNOO, BNSS, BNSeSe, BNTeTe, BNPoPo, and BNCOCO. ktoR(S1), 
ktoR(T1), and ktoR(T2) are the contributions from S1→S0 fluorescence including TADF, T1→S0 
phosphorescence, and T2→S0 phosphorescence to ktoR, respectively: ktoR=ktoR(S1)+ktoR(T1) 
+ktoR(T2); ktoR(S1)=kF(S1→S0)×[S1]/([S1]+[T1]+[T2]); ktoR(T1)=kPhos(T1S0)×[T1]/([S1]+[T1]+ 
[T2]); ktoR(T2)=kPhos(T2→S0)×[T2]/([S1]+[T1]+[T2]). ktoR for BNOO/BNSS/BNSeSe/BNTeTe/ 
BNPoPo/ BNCOCO is constant in the time domain longer than 100/100/10/10/1/10 ns. 


--- Page 27 ---
 
S27 
 
 
Supplementary Figure 3 | Calculated population, population ratio, and rate constants for 
reverse intersystem crossing (RISC) for BNOO, BNSS, BNSeSe, BNTeTe, BNPoPo, and 
BNCOCO. ktoRISC(T1) and ktoRISC(T2) are the contributions from T1→S1 and T2→S1 RISCs to 
ktoRISC, respectively: ktoRISC=ktoRISC(T1)+ktoRISC(T2); ktoRISC(T1)=kRISC(T1→S1)×[T1]/([T1]+[T2]); 
ktoRISC(T2) =kRISC(T2→S1)×[T2]/([T1]+[T2]). ktoRISC for BNOO/BNSS/BNSeSe/BNTeTe/ 
BNPoPo/BNCOCO is constant in the time domain longer than 1/10/10/10/1/10 ns. 


--- Page 28 ---
 
S28 
Supplementary Method 5 
Calculation of emission spectra 
The emission spectra were calculated using the vertical gradient method implemented in 
ORCA 5.0.3 program package6-8 for the S0 geometry optimised at the PBE0/Def2-SVP level. 
Frequency analysis was also performed at the same level. The input file for the geometry 
optimisation, frequency analysis, and spectrum calculations were as follows. 
 
# Geometry optimisation and frequency analysis 
!PBE0 Def2-SVP Opt Freq 
* XYZFILE 0 1 S0_init.xyz 
 
# Spectrum simulation 
!PBE0 DEF2-SVP TIGHTSCF ESD(FLUOR) 
%TDDFT NROOTS 1 
IROOT 1 
END 
%ESD 
GSHESSIAN "S0.hess" # Use the Hessian calculated at the optimised S0-state geometry 
DOHT TRUE 
HESSFLAG VG # Use the vertical gradient method 
LINES Gauss # Use the Gaussian distribution function 
INLINEW 679.40552 
SPECRES 5.0 # Spectral resolution 
UNIT NM 
END 
* XYZFILE 0 1 S0.xyz # Use the optimised S0-state geometry 
 
 


--- Page 29 ---
 
S29 
 
Supplementary Figure 4 | Experimental and calculated emission spectra for BNOO. The 
grey curve shows the experimental spectrum for 1 wt% BNOO-doped DMIC-TRZ film; the 
coloured curves show the calculated emission spectra. The Gaussian distribution functions 
with FWHM values of 1300 (green), 1400 (pink), 1500 (orange), 1600 (red), 1700 (blue), and 
1800 cm−1 (black) were examined to reproduce the experimental spectrum. The best fit was 
obtained for the FWHM values of 1600 cm−1 (red). 
 
 


--- Page 30 ---
 
S30 
 
Supplementary Figure 5 | Experimental and calculated emission spectra for BNSS. The 
grey curve shows the experimental spectrum for 1 wt% BNSS-doped DMIC-TRZ film; the 
coloured curves show the calculated emission spectra. The Gaussian distribution functions 
with FWHM values of 1300 (green), 1400 (pink), 1500 (orange), 1600 (red), 1700 (blue), and 
1800 cm−1 (black) were examined to reproduce the experimental spectrum. The best fit was 
obtained for the FWHM values of 1700 cm−1 (red). 
 
 


--- Page 31 ---
 
S31 
 
Supplementary Figure 6 | Experimental and calculated emission spectra for BNSeSe. The 
grey curve shows the experimental spectrum for 1 wt% BNSeSe-doped DMIC-TRZ film; the 
coloured curves show the calculated emission spectra. The Gaussian distribution functions 
with FWHM values of 1300 (green), 1400 (pink), 1500 (orange), 1600 (red), 1700 (blue), and 
1800 cm−1 (black) were examined to reproduce the experimental spectrum. The best fit was 
obtained for the FWHM values of 1600 cm−1 (red). 
 
 
 


--- Page 32 ---
 
S32 
 
Supplementary Figure 7 | Calculated emission spectra for BNTeTe, BNPoPo, and 
BNCOCO. The Gaussian distribution functions with FWHM values of 1600 (red) and 1700 
cm−1 (blue) were applied. 
 
 


--- Page 33 ---
 
S33 
Supplementary Method 6 
Molecular design for minimizing T1-T2 energy gap 
Supplementary Figure 8a shows a case when T1 and T2 consist of HOMO → LUMO and 
HOMO−1 → LUMO transitions, respectively. The T1 and T2 energies can be expressed as11 
𝐸(T1) = 2ℎH−1H−1 + ℎHH + ℎLL + 𝐽H−1H−1 + 2𝐽H−1H + 2𝐽H−1L + 𝐽HL − 𝐾H−1H − 𝐾HL 
(D1) 
𝐸(T2) = ℎH−1H−1 + 2ℎHH + ℎLL + 𝐽HH + 2𝐽H−1H + 𝐽H−1L + 2𝐽HL − 𝐾H−1H − 𝐾HL 
(D2) 
where h denotes the core integral (kinetic and potential energies), J denotes the Coulomb 
integral, and K denotes the exchange integral. Then, ∆𝐸(T2 − T1) is written as 
∆𝐸(T2 − T1) = (ℎHH − ℎH−1H−1) + (𝐽HH − 𝐽H−1H−1) + (𝐽HL − 𝐽H−1L) 
(D3) 
A simple approach of minimizing ΔE(T1→T2) is to decrease the first term hHH − hH−1H−1, 
which is possible by expanding the HOMO and HOMO−1 distributions (the second and third 
terms of ΔE(T1→T2), expressed in terms of the Coulomb integrals, are not easy to control). 
Supplementary Figure 8b shows a case where T1 and T2 consist of the HOMO → LUMO and 
HOMO → LUMO+1 transitions, respectively. The T1 and T2 energies can be expressed as 
𝐸(T1) = ℎHH + ℎLL + 𝐽HL − 𝐾HL 
(D4) 
𝐸(T2) = ℎHH + ℎL+1L+1 + 𝐽HL+1 − 𝐾HL+1 
(D5) 
Hence, ∆𝐸(T2 − T1) is written as 
∆𝐸(T2 − T1) = (ℎL+1L+1 − ℎLL) + (𝐽HL+1 − 𝐽HL) − (𝐾HL+1 − 𝐾HL) 
(D6) 
For the same reason, expanding the LUMO+1 and LUMO distributions results in an effective 
approach to minimize ΔE(T1→T2) by decreasing hL+1L+1 − hLL. Regardless of whether T2 is 
described as the HOMO−1 → LUMO or HOMO → LUMO+1 transition, expanding 
molecular orbitals relevant for T1 and T2 is a simple way to decrease ΔE(T1→T2) and 
accelerate the T2-mediated RISC process. 
 
Supplementary Figure 8 | S0, T1, and T2 determinants and ΔE(T1→T2). (a) A case when 
T1 and T2 consist of HOMO → LUMO (H → L) and HOMO−1 → LUMO (H−1 → L) 
transitions, respectively. (b) A case when T1 and T2 consist of HOMO → LUMO (H → L) 
and HOMO → LUMO+1 (H → L+1) transitions, respectively. 


--- Page 34 ---
 
S34 
Supplementary Table 15 | Minimization of ΔE(T1→T2). These examples demonstrate that 
expanding molecular orbitals relevant for T1 and T2 is an effective way to decrease ΔE(T1→
T2), resulting in acceleration of the T2-mediated RISC process. 
MR-TADF molecules with ΔE(T2–T1) 
References 
(in the main text) 
   
 
0.147 eV           >           0.098 eV 
12 
(45) 
    
 
0.28 eV    >    0.12 eV 
13 
(46) 
    
 
0.46 eV          >           0.33 eV 
14 
(47) 
    
 
0.21 eV   >     0.07 eV 
15 
(48) 
    
 
0.58 eV   >     0.28 eV 
16 
(49) 
   
 
0.81 eV    >    0.04 eV 
17 
(50) 
 


--- Page 35 ---
 
S35 
Supplementary References 
 
1. Shizu, K. & Kaji, H. Theoretical determination of rate constants from excited states: 
application to benzophenone. J. Phys. Chem. A 125, 9000-9010 (2021). 
2. Shizu, K. & Kaji, H. Comprehensive understanding of multiple resonance thermally 
activated delayed fluorescence through quantum chemistry calculations. Commun. Chem. 
5, 53 (2022). 
3.   Frisch, M. J. et al. Gaussian 16 Rev. C.01.  (2016). 
4.   McMurchie, L. E. & Davidson, E. R. One- and two-electron integrals over cartesian 
gaussian functions. J. Comput. Phys. 26, 218-231 (1978). 
5.   Yersin, H., Rausch, A. F., Czerwieniec, R., Hofbeck, T. & Fischer, T. The triplet state of 
organo-transition metal compounds. Triplet harvesting and singlet harvesting for efficient 
OLEDs. Coord. Chem. Rev. 255, 2622-2652 (2011). 
6. Neese, F. The ORCA program system. Wiley Interdiscip. Rev.: Comput. Mol. Sci. 2, 73-78 
(2012). 
7. Neese, F. Software update: the ORCA program system, version 4.0. WIREs Comput. Mol. 
Sci. 8, e1327 (2017). 
8. Neese, F., Wennmohs, F., Becker, U. & Riplinger, C. The ORCA quantum chemistry 
program package. J. Chem. Phys. 152, 224108 (2020). 
9. Hu, Y. X., Miao, J., Hua, T., Huang, Z., Qi, Y., Zou, Y., Qiu, Y., Xia, H., Liu, H., Cao, X. 
& Yang, C. Efficient selenium-integrated TADF OLEDs with reduced roll-off. Nat. 
Photon. 16, 803–810 (2022) 
10. TURBOMOLE V7.4.1 2019, a development of University of Karlsruhe and 
Forschungszentrum Karlsruhe GmbH, 1989–2007, TURBOMOLE GmbH, since 2007. 
11. Szabo, A. & Ostlund, N. S. Modern quantum chemistry: Introduction to advanced 
electronic structure theory 87-89 (Dover Publications, New York, 1996). 
12. Oda, S. et al. One-shot synthesis of expanded heterohelicene exhibiting narrowband 
thermally activated delayed fluorescence. J. Am. Chem. Soc. 144, 106-112 (2022). 
13. Dos Santos, J. M. et al. An s-shaped double helicene showing both multi-resonance 
thermally activated delayed fluorescence and circularly polarized luminescence. J. Mater. 
Chem. C 10, 4861-4870 (2022). 
14. Hall, D. et al. Diindolocarbazole – achieving multiresonant thermally activated delayed 
fluorescence without the need for acceptor units. Mater. Horiz. 9, 1068-1080 (2022). 
15. Jin, J. et al. Integrating asymmetric O−B−N unit in multi-resonance thermally activated 
delayed fluorescence emitters towards high-performance deep-blue organic light-emitting 
diodes. Angew. Chem. Int. Ed. 62, e202218947 (2023). 
16. Wu, S. et al. Merging boron and carbonyl based mr-tadf emitter designs to achieve high 
performance pure blue oleds. Angew. Chem. Int. Ed. 62, e202305182 (2023). 
17. Keruckiene, R. et al. Is a small singlet–triplet energy gap a guarantee of TADF 
performance in MR-TADF compounds? Impact of the triplet manifold energy splitting. J. 
Mater. Chem. C 12, 3450-3464 (2024).